\documentclass[a4paper,11pt,fleqn]{article}%fleqn pour aligner les equations à gauche
\usepackage{Cours}


\newcommand{\chnum}{Chapitre 3}
\newcommand{\chtitle}{Titrage par avec suivi conductimétrique - \emph{Correction}}
\newcommand{\dtype}{TP}


\begin{document}
\maketitle

\section{Détermination de l'équivalence lors d'un titrage par suivi conductimétrique}
	\subsection{Analyser et Raisonner}

\begin{enumerate}
	\item Matériel pour la dilution : fiole jaugée de \SI{1}{L} ; pipette jaugée de \SI{10}{mL} ; propipette ; solution $S_0$ ; eau distillée.
	\item Pour réaliser un suivi conductimétrique il faut que l'équation support de la réaction implique des ion. C'est le cas ici : \ce{H3O+(aq) + HO-(aq) -> 2H2O(l)}\\
	Il faut aussi une solution assez diluée ce qui devrait être le cas avec la solution $S_1$ qui est diluée 100 fois.
	\item Protocole expérimental :
	\begin{itemize*}
		\item Prélever \SI{10}{mL} de solution $S_1$ à l'aide d'une pipette jaugée et les placer dans un grand bécher.
		\item Prélever \SI{200}{mL} d'eau distillée à l'aide d'une éprouvette graduée et les ajouter dans le bécher.
		\item Rincer la burette graduée avec la solution d'acide chlorhydrique à \SI{3.00E-2}{\mole\per\liter} puis remplir la burette.
		\item Placer un barreau aimanté dans le grand bécher
		\item Placer le bécher sur un agitateur magnétique et mettre l'ensemble sous la burette graduée.
		\item Placer la sonde du conductimètre étalonné dans la solution et commencer l'agitation.
		\item Ajouter la solution titrante mL par mL jusqu'à \SI{20}{mL} en notant les valeurs de conductivité à chaque fois.
		\item Rentrer les valeurs dans LatisPro
		\item Utiliser la fonction "Tracer droite" pour tracer les droites tangentes à la courbe avant et après la rupture de courbe.
		\item Relever le volume équivalent $V_E$ correspondant à l'abscisse de l'intersection des deux droites.
	\end{itemize*}
\end{enumerate}

\subsection{Réaliser}

\begin{enumerate}[resume]
	\item (manipulation)
	\item On obtient $V_E \approx$ \SI{9.8}{mL}. 
	\item On calcule $c_1$ à partir des données à l'équivalence. $$n_1=n_E$$ $$\leftrightarrow c_1 \cdot V_1 = c_A \cdot V_E$$ $$\leftrightarrow c_1 = \dfrac{c_A \cdot V_E}{V_1}$$ 
	\fbox{AN : } $ c_1 = \dfrac{\num{3.00E-2}\times \num{9.8E-3}}{\num{10E-3}}= \SI{2.94E-2}{\mole\per\liter}$
	\item La solution $S_0$ est 100 fois plus concentrée que la solution $S_1$ ce qui donne une concentration $c_0 = 100 \times c_1$\\
	\fbox{AN : } $c_0 = 100 \times \num{2.94E-2} = \SI{2.94}{\mole\per\liter}$
\end{enumerate}

\subsection{Valider}

 \begin{enumerate}[resume]
	\item Le fabricant donne $C_m $ = \SI{120}{\gram\per\liter}. $$c_{0, fabricant} = \dfrac{C_m}{M}$$
	\fbox{AN : } $c_{0, fabricant} = \dfrac{\num{120}}{\num{40}} = \SI{3}{\mole\per\liter}$
	\item En observant le matériel on note les incertitudes de mesure suivantes : $u(V_1)$ = \SI{0.040}{mL} (pipette jaugée de \SI{10}{mL}) et $u(V_A,E)$ = \SI{0.1}{mL} (burette graduée de \SI{25}{mL}). Pour déterminer l'incertitude $u(c_0)$ on modifie la formule donnée, soit :
	$$u(c_0) = c_0 \cdot \sqrt{\left(\frac{u(c_A)}{c_A}\right)^2+\left(\frac{u(V_{A,E})}{V_{A,E}}\right)^2+\left(\frac{u(V_1)}{V_1}\right)^2}$$
	\fbox{AN : } $u(c_0) =\num{2.94} \cdot \sqrt{\left(\dfrac{\num{0.001}}{\num{3.00E-2}}\right)^2+\left(\dfrac{\num{0.1E-3}}{\num{9.8E-3}}\right)^2+\left(\dfrac{\num{0.040E-3}}{\num{10E-3}}\right)^2} = \SI{3.94E-2}{\mole\per\liter}$
	\item L'écart entre la valeur expérimentale et la valeur annoncée $\Delta c_0$ donne : $$\Delta c_0 = c_{0,fabricant} - c_0$$
	\fbox{AN : } $\Delta c_0 = \num{3}-\num{2.94} = \SI{6E-2}{\mole\per\liter}$\\
	On constate que la différence dépasse l'intervalle de confiance calculé précédemment : Le titrage n'est pas assez précis.
\end{enumerate}
\section*{Comprendre la courbe de titrage}
\begin{enumerate}[resume]
	\item 
\end{enumerate}
\noindent \begin{tabular}{|>{\centering}m{3.4cm}|>{\centering}m{3.4cm}|>{\centering}m{3.4cm}|>{\centering}m{3.4cm}|>{\centering\arraybackslash}m{3.4cm}|} \hline
	\textbf{Temps} & \ce{HO-} & \ce{Na+} & \ce{H3O+} & \ce{Cl-}\\ \hline
	Avant l'équivalence & $\downarrow$ & $=$ & $0$ & $\uparrow$\\ \hline
	A l'équivalence & $0$ & $=$ & $0$ &$\uparrow$\\ \hline
	Après l'équivalence & $0$ & $=$ & $\uparrow$ & $\uparrow$ \\ \hline
\end{tabular}
\begin{enumerate}[resume]
	\item La conductivité peut se calculer par la formule : $$\sigma = \sum{\lambda_i \cdot [X_i]}$$
	\item Avant l'équivalence, les ions en présence dans le milieu sont \ce{HO-}, \ce{Na+} et \ce{Cl-}, la conductivité va doucement diminuer avec la consommation des ion \ce{HO-} et ce malgré l'augmentation des ions \ce{Cl-} dans le milieu car $\lambda (\ce{HO_}) >\lambda (\ce{Cl-})$.\\
	Après l'équivalence, les ions en présence dans le milieu seront \ce{Na+}, \ce{H3O+} et \ce{Cl-}. La conductivité va augmenter avec l'ajout continu d'ions \ce{H3O+} et \ce{Cl-} dans le milieu.
\end{enumerate}
\end{document}